% =============================================================================
% 02_motivation.tex -- Introduction and Motivation (~5-6 frames)
% =============================================================================

\section{Motivation}

% ---- Frame 1 ----
\begin{frame}{What Is a Brain Digital Twin?}

  A \alert{Brain Digital Twin} is a subject-specific computational replica
  that can simulate mechanical processes in an individual brain
  \cite{Goriely2015,Meaney2014}.

  \bigskip

  \begin{columns}[T]
    \column{0.55\textwidth}
    \begin{itemize}
      \item Personalized at every layer:
      \begin{itemize}
        \item Geometry from \textbf{this patient's} T1w
        \item Fiber architecture from \textbf{this patient's} dMRI
        \item Functional context from \textbf{this patient's} fMRI
      \end{itemize}
      \item Continuously informed by imaging data
      \item Predicts response to perturbations
    \end{itemize}

    \column{0.4\textwidth}
    % FIGURE SLOT: schematic of digital twin concept
    % \includegraphics[width=\textwidth]{dt_concept_schematic.png}
    \begin{center}
      \fbox{\parbox{0.9\textwidth}{\centering\small
        \textit{Figure:} Digital twin concept\\
        (physical brain $\leftrightarrow$ virtual model)
      }}
    \end{center}
  \end{columns}

\end{frame}

% ---- Frame 2 ----
\begin{frame}{Why Gliomas?}

  Diffuse gliomas present unique challenges for brain mechanics:

  \bigskip

  \begin{itemize}
    \item \textbf{Clinical urgency:} most common primary brain tumors;
      median survival for GBM $\sim$15 months
    \item \textbf{Mass effect:} tumor growth displaces and compresses
      surrounding tissue \cite{Zarodniuk2026}
    \item \textbf{Surgical brain shift:} resection creates cavity,
      CSF dynamics change, tissue deforms
    \item \textbf{Longitudinal changes:} pre-op $\to$ post-op provides
      \emph{empirical deformation data}
    \item \textbf{Pathological anatomy:} standard pipelines break ---
      segmentation, registration, boundary conditions all need adaptation
  \end{itemize}

  \bigskip
  \alert{$\Rightarrow$} Glioma patients are the \emph{hardest} and most
  \emph{clinically relevant} test case for brain digital twins.

\end{frame}

% ---- Frame 3 ----
\begin{frame}{The Finse Workshop Context}

  \textbf{CAS -- Mathematical Challenges in Brain Mechanics}\\
  Finse, March 2026

  \bigskip

  \begin{itemize}
    \item Applied mathematicians $\times$ neuroscientists $\times$
      clinicians
    \item Focus on the \emph{mathematical} difficulties that arise when
      building subject-specific brain models
    \item This talk: a \emph{concrete worked example} grounded in real
      patient data
  \end{itemize}

  \bigskip

  \textbf{Companion project:} Single-subject healthy volunteer
  (\texttt{brain-digital-twin-finse2026})\\
  $\to$ establishes the pipeline; \textbf{this project} extends it to
  pathology and longitudinal change.

\end{frame}

% ---- Frame 4 ----
\begin{frame}{The Governing Equation}

  The mechanical forward problem:
  \[
    \nabla \cdot \bfsigma(\bfu, \bftheta) + \bfb
    = \rho\,\ddot{\bfu}
  \]

  \begin{itemize}
    \item $\bfu(\bfx,t)$ -- displacement field
    \item $\bfsigma$ -- Cauchy stress (from constitutive law)
    \item $\bftheta(\bfx)$ -- spatially varying material parameters
      (stiffness, anisotropy, viscosity)
    \item $\bfb$ -- body forces; $\rho$ -- tissue density
  \end{itemize}

  \bigskip

  MRI informs: the domain $\Omega$ (geometry), the parameter field
  $\bftheta$ (material), and -- partially -- the boundary conditions.

  \bigskip

  \alert{Core difficulty:} non-uniqueness. Many $(\bftheta, \text{BC})$
  combinations explain the same images.

\end{frame}

% ---- Frame 5 ----
\begin{frame}{This Talk: Roadmap}

  \begin{enumerate}
    \item Neuro-oncological challenges
    \item Multiparametric MRI opportunities
    \item Data and metadata (BTC dataset, BIDS)
    \item Processing pipeline (9 steps)
    \item Longitudinal registration and deformation fields
    \item The digital twin concept and architecture
    \item Mathematical challenges
    \item Results preview and outlook
  \end{enumerate}

  \bigskip
  \small
  \textit{Supplemented by Google Slides with Freeview illustrations
  (static and dynamic).}

\end{frame}
